Un point $M$ situé sur une circonférence de rayon $R=\qty{1}{\meter}$ décrit un
mouvement dont l'équation horaire est~:
\[
    \theta(t)=\frac{\pi}{2}+2 \cdot t \qquad (\unit{\rad})
\]
$\theta$~: abscisse angulaire à l'instant $t$ et $\theta_{0}$ abscisse
angulaire à la date $t=\qty{0}{\second}$.

Sur un schéma et à la date $t=\qty{2}{\second}$, représenter~:
\begin{enumerate}
    \item La position angulaire du point $M$.
    \item Le vecteur vitesse du point $M$.
\end{enumerate}
Échelle~: $\qty{4}{\cm} \to \qty{1}{\meter}$ et $\qty{2}{\cm} \to \qty{1}{\v}$.

